\documentclass[10pt]{article}

%%% Doc layout
\usepackage{fullpage} 
\usepackage{booktabs}       % professional-quality tables
\usepackage{microtype}      % microtypography
\usepackage{parskip}
\usepackage{times}
\usepackage{graphicx}

%% Hyperlinks always black, no weird boxes
\usepackage[hyphens]{url}
\usepackage[colorlinks=true,allcolors=black,pdfborder={0 0 0}]{hyperref}

%%% Math typesetting
\usepackage{amsmath,amssymb}

%%% Write out problem statements in blue, solutions in black
\usepackage{xcolor}
\newcommand{\officialdirections}[1]{{\color{purple} #1}}

%%% Avoid automatic section numbers (we'll provide our own)
\setcounter{secnumdepth}{0}

%% --------------
%% Header
%% --------------
\usepackage{fancyhdr}
\fancyhf{}
\setlength{\headheight}{15pt}
\fancyhead[C]{\ifnum\value{page}=1 Tufts CS 145 BDL - 2026s - HW2 Submission \else \fi}
\fancyfoot[C]{\thepage} % page number
\renewcommand\headrulewidth{0pt}
\pagestyle{fancy}

%% --------------
%% Begin Document
%% --------------
\begin{document}

~\\ %% add vertical space
{\Large{\bf Student Name: TODO}}

~\\ %% add vertical space
{\bf Collaboration Statement:}

Total hours spent: TODO

I consulted the following human, textbook, or AI resources:
\begin{itemize}
\item TODO
\item TODO
\item $\ldots$	
\end{itemize}

By turning this document in, I attest that I have followed the 
\href{https://www.cs.tufts.edu/cs/145/2026s/index.html#collaboration}{[CS 145 Collaboration Policy]}. All solutions represent my own work. No solution text in this document was directly provided by other humans or artificial agents. 

\tableofcontents

\newpage

\subsection{1a: Plot of potential energy vs. iteration for HMC with 0 hiddens}

\renewcommand{\figurename}{Fig.}
\renewcommand{\thefigure}{1a}
 \begin{figure}[!h]
     \centering
     \includegraphics[height=5cm,width=0.9\textwidth,keepaspectratio]{example-image-a.pdf}
     \label{fig:1a}
\caption{TODO YOUR BRIEF CAPTION HERE. 
}%endcaption
 \end{figure}


\subsection{1b: Plot of posterior predictive for HMC with 0 hiddens}

\renewcommand{\thefigure}{1b}
 \begin{figure}[!h]
     \centering
     \includegraphics[height=12.5cm,width=0.9\textwidth,keepaspectratio]{example-image-b.pdf}
     \label{fig:1b}
\caption{TODO YOUR BRIEF CAPTION HERE. 
}%endcaption
 \end{figure}

\newpage

\officialdirections{
\subsection{1c: Convergence assessment (short answer)}
Did your MCMC chains converge to the intended posterior? How can you tell?
}

TODO YOUR SOLUTION HERE

\officialdirections{
\subsection{1d: Tuning assessment (short answer)}
How did you set the \texttt{step\_size} and \texttt{n\_leapfrog\_steps} values? What other values did you try? How did you know when a setting was good enough? 
}

TODO YOUR SOLUTION HERE

\newpage

\subsection{2a: Plot of potential energy vs. iteration for Exact with 0 hiddens}

\renewcommand{\figurename}{Fig.}
\renewcommand{\thefigure}{2a}
 \begin{figure}[!h]
     \centering    \includegraphics[height=5cm,width=0.9\textwidth,keepaspectratio]{example-image-a.pdf}
     \label{fig:2a}
\caption{TODO YOUR BRIEF CAPTION HERE. 
}%endcaption
 \end{figure}


\subsection{2b: Plot of posterior predictive for Exact with 0 hiddens}

\renewcommand{\thefigure}{2b}
 \begin{figure}[!h]
     \centering
     \includegraphics[height=12.5cm,width=0.9\textwidth,keepaspectratio]{example-image-b.pdf}
     \label{fig:2b}
\caption{TODO YOUR BRIEF CAPTION HERE. 
}%endcaption
 \end{figure}



\newpage

\officialdirections{
\subsection{2c: HMC vs. exact (short answer)}
Does the sampled distribution from HMC match the exact posterior? How can you tell?
}

TODO YOUR SOLUTION HERE

\officialdirections{
\subsection{2d: Purpose of standardization (short answer)}
Why do we use a standardized version of the data? What's the value of standardizing the x? What's the value of standardizing the y?
}

TODO YOUR SOLUTION HERE


\newpage

\subsection{3a: Plot of potential energy vs. iteration for HMC with 64 hiddens}

\renewcommand{\figurename}{Fig.}
\renewcommand{\thefigure}{3a}
 \begin{figure}[!h]
     \centering
     \includegraphics[height=5cm,width=0.9\textwidth,keepaspectratio]{example-image-a.pdf}
     \label{fig:3a}
\caption{TODO YOUR BRIEF CAPTION HERE. 
}%endcaption
 \end{figure}


\subsection{3b: Plot of posterior predictive for Exact with 0 hiddens}

\renewcommand{\thefigure}{3b}
 \begin{figure}[!h]
     \centering
     \includegraphics[height=12.5cm,width=0.9\textwidth,keepaspectratio]{example-image-b.pdf}
     \label{fig:3b}
\caption{TODO YOUR BRIEF CAPTION HERE. 
}%endcaption
 \end{figure}



\newpage

\officialdirections{
\subsection{3c: Convergence assessment for 64 hiddens (short answer)}
Did each of your MCMC chains converge to the intended posterior? How can you tell?
}

TODO YOUR SOLUTION HERE

\officialdirections{
\subsection{3d: Tuning assessment for 64 hiddens (short answer)}
How did you set the \texttt{step\_size} and \texttt{n\_leapfrog\_steps} values? What other values did you try? How did you know when a setting was good enough? 
}

TODO YOUR SOLUTION HERE


\end{document}
